\chapter{Practices Leading to Meditation}
\label{ch:meditation}

\emph{How does a practitioner move from posture and breath towards meditation, and how is a first sitting taught?}
The syllabus closes the practical demonstrations with meditation and the practices that lead to it: recitation, body and breath awareness, and yoga nidra. This chapter gives the meditation of the common protocol in full and sets out the preparatory practices as far as the sources reached allow.

\section{What leads to meditation}

The Board syllabus names the practices that lead to meditation as dharana (concentration) and dhyana (meditation), with allied practices such as mudras, mantra and japa, and a suitable place for sitting. Everything earlier in the session serves this: the loosening practices and asanas make the body steady and comfortable, so that it can sit still; pranayama makes the breath slow, and with it the mind; recitation gives the mind a single sound to rest on.

One teaching sequence used in the workshop runs from mindfulness and pranayama, through stable asanas, to concentration and visualisation. It is one useful way to prepare, not an order set by the Board; the classical order is Patanjali's eight limbs, from yama and niyama to dharana, dhyana and samadhi (Yoga Sutra 2.29).

\section{Dhyana in jnana mudra}

The common protocol gives this method:
\begin{enumerate}
\item Sit in any meditative posture with the spine comfortably upright.
\item Hold jnana mudra: the tip of the thumb touching the tip of the index finger, the other three fingers straight, relaxed and side by side, the palms up on the thighs.
\item Let the arms and shoulders be loose. Close the eyes and lift the face slightly.
\item Do not force concentration. Keep a mild attention between the eyebrows and stay aware of the breath.
\item Let the thoughts settle, gradually, towards one quiet thought.
\item Stay as long as is comfortable.
\end{enumerate}
Beginners may have soft music in the background.

\section{Teaching a first sitting}

Meditation is the heart of the practice, but a first session should be short. Say plainly that thoughts will come, and that noticing them and returning to the breath is part of the practice, not a failure. A student who strains to empty the mind becomes tense; the instruction is ``mild attention'', not ``stop thinking''. A student who cannot sit on the floor sits on a chair with the feet flat and the spine upright.

\section{Recitation, awareness and yoga nidra}

The syllabus also lists the recitation of pranava (om) and soham, body awareness, breath awareness, and yoga nidra. The common protocol does not describe them; the methods below follow Swami Satyananda Saraswati of the Bihar School of Yoga. The meaning of the pranava and of japa is in the chapter on prayer (see \cref{ch:prayer}).

\subsection*{Pranava (om) recitation}
Sit steady with the spine upright and the eyes closed. Chant AUM slowly on the out-breath, and rest in the silence after it. Repeat for a few rounds, then continue mentally. Keep the meaning in mind: the Yoga Sutra joins the repetition of the pranava to feeling its meaning (1.27 to 1.28).

\subsection*{Soham recitation}
This is \emph{ajapa japa}, the mantra that goes with the breath. Sit steady and relax for a few minutes, and watch the natural breath. Then hear ``so'' with each in-breath and ``ham'' with each out-breath, without forcing or changing the breath. Some tantric teachings reverse the two syllables; the form above is the Bihar School's.

\subsection*{Body awareness}
Lie in shavasana. Move the attention quickly from one body part to the next in a fixed order: start with the right hand thumb, go through the rest of the right side, then the left side, then the back, the front and the head. Name each part mentally and feel it, without moving. Satyananda calls this the rotation of consciousness.

\subsection*{Breath awareness}
Watch the natural breath without altering or forcing it. Only observe it.

\subsection*{Yoga nidra}
Yoga nidra is a guided practice done lying in shavasana. Satyananda's form has eight stages: settling in; a short positive resolve (\emph{sankalpa}), repeated mentally three times; rotation of consciousness; breath awareness; pairs of opposite feelings such as heaviness and lightness; visualisation; the resolve repeated; and coming out slowly. The student stays awake and aware throughout. Satyananda presented it as ancient; scholars note that the multi-stage form is modern and not found in the old texts.

\section{Closing the session}

The session ends with a sankalpa and the peace chant, as described in the chapter on prayer.

\begin{recall}
\begin{itemize}
\item Syllabus: dharana and dhyana, with mudra, mantra, japa and a suitable place.
\item Protocol dhyana: meditative seat, jnana mudra, eyes closed, mild attention between the eyebrows, aware of the breath.
\item Do not force concentration; thoughts are expected.
\item Close with a sankalpa and the peace chant.
\end{itemize}
\end{recall}

\subsection*{Worked examples}
\begin{itemize}
\item \textbf{A student says they ``failed'' because their mind kept wandering.} Explain that noticing the wandering and coming back to the breath is the practice itself. Keep the next sitting shorter.
\item \textbf{A student with stiff knees cannot sit cross-legged.} Seat them on a chair, feet flat, spine upright, hands in jnana mudra on the thighs.
\end{itemize}

\subsection*{Check yourself}
\begin{enumerate}
\item Which practices does the Board syllabus name as leading to meditation?
\item Where does the protocol place the mild attention in meditation?
\item Which hand gesture is used?
\item Why are asanas practised before sitting for meditation?
\item How does the session end?
\end{enumerate}

\subsection*{Answers}
\begin{enumerate}
\item Dharana and dhyana, with allied practices such as mudras, mantra and japa, and a suitable environment.
\item Between the eyebrows, with awareness of the breath.
\item Jnana mudra.
\item To make the body stable and comfortable for sitting.
\item With a sankalpa and the peace chant.
\end{enumerate}

\paragraph{Sources.} Common Yoga Protocol, Ministry of AYUSH (dhyana); the Board syllabus for Level 1; the official guidebook for Level I (2016, older scheme); the Yoga Sutra 1.27 to 1.28 (pranava japa); Swami Satyananda Saraswati, \emph{Yoga Nidra} (Yoga Publications Trust, Munger), with the stages as listed in Wikipedia, ``Yoga nidra'', and the rotation and sankalpa as given in Sanskriti magazine, ``Yoga Nidra Techniques''; Satyananda on soham as quoted from \emph{Yoga} magazine (Bihar School of Yoga); on the modern origin of the multi-stage yoga nidra, J. Birch and J. Hargreaves, as summarised in Wikipedia, ``Yoga nidra''.
